% !TeX program = pdflatex
% DLR-style cover letter for Journal of Fluid Mechanics submission.
% Compile with: pdflatex jfm_cover_letter.tex
%
% Put dlr_logo.png in the same folder as this .tex file.

\documentclass[11pt,a4paper]{article}

\usepackage[a4paper,left=25mm,right=25mm,top=22mm,bottom=22mm]{geometry}
\usepackage{graphicx}
\usepackage[absolute,overlay]{textpos}
\usepackage{xcolor}
\usepackage{helvet}
\usepackage[T1]{fontenc}
\usepackage[utf8]{inputenc}
\usepackage[english]{babel}
\usepackage{setspace}
\usepackage{parskip}
\usepackage{ragged2e}
\usepackage{microtype}

\renewcommand{\familydefault}{\sfdefault}
\pagestyle{empty}
\setlength{\parindent}{0pt}
\setstretch{1.02}

% -----------------------------------------------------------------------------
% Editable letter data
% -----------------------------------------------------------------------------
\newcommand{\Institute}{Institut für Aerodynamik und Strömungstechnik}
\newcommand{\DLRAddress}{Lilienthalplatz 7, 38108 Braunschweig}

\newcommand{\Recipient}{%
The Editor\\
Journal of Fluid Mechanics\\
Cambridge University Press%
}

\newcommand{\ContactPerson}{Dr. Sparsh Sharma}
\newcommand{\Phone}{0531 295--1530}
\newcommand{\Email}{sparsh.sharma@dlr.de}
\newcommand{\LetterDate}{2 July 2026}

\newcommand{\Subject}{Submission of an original manuscript to the Journal of Fluid Mechanics}

% Use your final manuscript title here.
\newcommand{\ManuscriptTitle}{Strain-coupled one-dimensional turbulence for rapid distortion}

\newcommand{\ClosingName}{Dr. Sparsh Sharma}

% Logo file
\newcommand{\LogoFile}{dlrlogo.png}

% Optional: insert scanned signature image here if you have one.
% Leave blank if not available.
\newcommand{\SignatureImage}{} % e.g. signature.png

% -----------------------------------------------------------------------------
% Small helpers
% -----------------------------------------------------------------------------
\newcommand{\SmallGrey}[1]{\textcolor{black!55}{\scriptsize #1}}
\newcommand{\ContactLabel}[1]{\textcolor{black!55}{\scriptsize #1}}

\begin{document}

% -----------------------------------------------------------------------------
% Header
% -----------------------------------------------------------------------------
\begin{textblock*}{95mm}(25mm,28mm)
  {\small \Institute}
\end{textblock*}

\begin{textblock*}{75mm}(126mm,20mm)
  \raggedleft
  \includegraphics[width=58mm]{\LogoFile}
\end{textblock*}

% -----------------------------------------------------------------------------
% Recipient address block and contact block
% -----------------------------------------------------------------------------
\begin{textblock*}{86mm}(25mm,66mm)
  {\scriptsize
  DLR e. V.\hspace{3mm} \Institute\\[-0.2mm]
  \hspace*{21mm}\DLRAddress}
  \vspace{7mm}

  {\normalsize \Recipient}
\end{textblock*}

\begin{textblock*}{65mm}(126mm,66mm)
  \raggedleft
  \begin{tabular}{@{}r@{\hspace{7mm}}l@{}}
    \ContactLabel{Reference}       & \\
    \ContactLabel{Manuscript type} & JFM Papers \\
    \\[-1mm]
    \ContactLabel{Contact person}  & \ContactPerson \\
    \\[-1mm]
    \ContactLabel{Telephone}       & \Phone \\
    \\[-1mm]
    \ContactLabel{E-mail}          & \Email \\
    \\[2mm]
                                    & \LetterDate
  \end{tabular}
\end{textblock*}

% -----------------------------------------------------------------------------
% Body
% -----------------------------------------------------------------------------
\vspace*{118mm}

{\bfseries \Subject}

\vspace{11mm}

Dear Editor,

\vspace{5mm}

I am pleased to submit the manuscript entitled
``\ManuscriptTitle'' for consideration as a regular article in the
\textit{Journal of Fluid Mechanics}.

The manuscript develops a strain-coupled formulation of one-dimensional
turbulence for flows undergoing rapid distortion by an imposed mean strain.
The central objective is to extend ODT beyond its standard unstrained form
by incorporating the production, rapid pressure--strain redistribution and
kinematic wavenumber compression associated with homogeneous mean strain,
while preserving the structural ingredients of the model. The formulation is
verified against rapid-distortion theory and is intended as a reduced-order,
scale-resolved route to predicting distorted anisotropic turbulence spectra,
with leading-edge-noise applications as the motivating context.

We believe the manuscript is suitable for JFM because it addresses a
fundamental modelling problem in turbulence: how rapid mean-strain
distortion, pressure--strain redistribution and nonlinear cascade dynamics
can be represented consistently in a reduced-order framework. The paper
therefore combines analytical model construction, verification against
rapid-distortion limits, and comparison with canonical strained-turbulence
data.

The manuscript is original, has not been published previously, and is not
under consideration elsewhere. All authors have approved the submission. To
the best of my knowledge, there are no conflicts of interest that would
affect the review of this work.

\vspace{5mm}

Thank you for considering our manuscript for publication in the
\textit{Journal of Fluid Mechanics}.

\vspace{10mm}

Yours sincerely,

\vspace{13mm}

% Signature placeholder. If you have a signature image, set \SignatureImage above.
\if\relax\detokenize\expandafter{\SignatureImage}\relax
  \vspace{8mm}
\else
  \includegraphics[height=15mm]{\SignatureImage}\\[-1mm]
\fi

\ClosingName

% -----------------------------------------------------------------------------
% Footer
% -----------------------------------------------------------------------------
\begin{textblock*}{65mm}(25mm,260mm)
  \scriptsize
  Deutsches Zentrum für Luft- und Raumfahrt e.\,V.\\
  Institut für Aerodynamik und Strömungstechnik\\
  Lilienthalplatz 7\\
  38108 Braunschweig
\end{textblock*}

% \begin{textblock*}{55mm}(139mm,260mm)
%   \scriptsize
%   Lilienthalplatz 7\\
%   38108 Braunschweig\\
%   Telefon 0531 295--0\\
%   Internet DLR.de
% \end{textblock*}

\end{document}